%% Springer Nature journal template -- sn-jnl (Version 3.1, December 2024)
%% Reference style: Math and Physical Sciences Numbered
\documentclass[pdflatex,sn-mathphys-num]{sn-jnl}

%%%% Packages
\usepackage{graphicx}
\usepackage{multirow}
\usepackage{amsmath,amssymb,amsfonts}
\usepackage{amsthm}
\usepackage{booktabs}
\usepackage{xcolor}
\usepackage[title]{appendix}
\usepackage{textcomp}

%%%% Theorem styles
\theoremstyle{thmstyleone}
\newtheorem{theorem}{Theorem}
\newtheorem{proposition}[theorem]{Proposition}

\theoremstyle{thmstyletwo}
\newtheorem{example}{Example}
\newtheorem{remark}{Remark}

\theoremstyle{thmstylethree}
\newtheorem{definition}{Definition}

\raggedbottom

\begin{document}

\title[Feature-Driven Meta-Learning for Decomposition Strategy Selection]{%
  When Does Seasonal Decomposition Help?\\
  A Feature-Driven Meta-Learning Study for Time Series Forecasting}

\author*[1]{\fnm{Diogo} \sur{Costa}}\email{diogocosta1104@gmail.com}

\affil*[1]{%
  \orgdiv{Faculty of Engineering},
  \orgname{University of Porto},
  \orgaddress{%
    \street{Rua Dr.\ Roberto Frias},
    \city{Porto},
    \postcode{4200-465},
    \country{Portugal}}}

\abstract{%
Seasonal decomposition is one of the oldest preprocessing tools in time series
forecasting, yet practitioners receive little guidance on when it actually
helps. We study this question empirically on 3{,}000 monthly and quarterly
series from the M3 and M4 competitions. Each series is described by six
interpretable features that capture non-normality, nonlinearity, spectral
entropy, evolving seasonality, structural breaks, and conditional
heteroscedasticity. We compare three strategies across ten automatically tuned
models from statistical, machine learning, neural, and transformer families:
forecasting the raw series, forecasting the seasonally adjusted series with a
naive seasonal add-back, and forecasting every STL component with a separate
model. We also test whether global models trained on feature-homogeneous
subsets of series outperform a single global model trained on everything. The
effect of decomposition turns out to be strongly family specific. Transformers
benefit consistently, with win rates near 68 percent. Statistical models are
harmed because they already model trend and seasonality internally. Machine
learning and neural models improve on most series but fail badly on a
minority, which pulls their average performance below the naive baseline.
Feature-based specialisation improves non-linear models and harms linear ones,
and pairing it with seasonal adjustment removes most catastrophic failures.
The results translate into simple, evidence-based rules for deciding when to
decompose.
}

\keywords{Time series forecasting, Seasonal decomposition, STL,
  Meta-learning, Feature engineering, AutoML, Global forecasting models}

\maketitle

% ---------------------------------------------------------------------------
\section{Introduction}\label{sec:intro}
% ---------------------------------------------------------------------------

Seasonal structure dominates many of the series that organisations forecast
every day, from retail sales to energy demand. A long-standing recipe handles
it in two steps: decompose the series, forecast the parts that are easier to
model, and add the seasonality back at the end. Seasonal and trend
decomposition using Loess, known as STL~\cite{cleveland1990stl}, has been the
workhorse of this recipe for three decades. It is robust, fast, implemented in
every major forecasting library, and easy to explain.

Whether the recipe still pays off is much less clear. Classical statistical
methods such as exponential smoothing and ARIMA model trend and seasonality
internally, so an external decomposition may simply duplicate work. Global
machine learning and deep learning models train on many series at once and
have no built-in notion of seasonality, so they might benefit from the cleaner
targets that decomposition produces. The winning method of the M4 competition
relied on internal deseasonalisation~\cite{smyl2020hybrid}, and recent neural
architectures embed decomposition blocks directly into the
network~\cite{wu2021autoformer, zeng2023dlinear}. Practice has clearly
absorbed the idea that decomposition matters. Evidence about when it helps,
for which models, and on which series remains thin.

This paper studies that question empirically. We describe each series by six
interpretable features computed before any forecasting model is fitted:
non-normality, nonlinearity, spectral entropy, evolving seasonality,
structural break strength, and conditional heteroscedasticity. We then
evaluate three decomposition strategies, from no decomposition at all to
forecasting every STL component with a separate model, across ten
automatically tuned models from statistical, machine learning, neural, and
transformer families. The testbed contains 3{,}000 monthly and quarterly
series sampled from the M3 and M4 competitions in a way that preserves their
feature distributions. A second experiment asks a complementary question: if
series are grouped by feature level, does a global model trained only on one
group beat a global model trained on everything?

The study makes four contributions.
\begin{enumerate}
  \item A systematic measurement of the effect of STL preprocessing across
    model families, showing that the effect is strongly family specific.
    Transformers gain, statistical models lose, and machine learning and
    neural models improve on most series while failing badly on a minority.
  \item An analysis of feature-conditioned training showing that
    specialisation helps non-linear models, harms linear ones, and is safest
    when combined with seasonal adjustment.
  \item Simple evidence-based rules that practitioners can apply directly,
    collected in Section~\ref{sec:summary} and discussed in
    Section~\ref{sec:discussion}.
  \item A documented failure mode of tercile-based feature bucketing when a
    feature has a degenerate distribution, together with the guard that
    prevents it.
\end{enumerate}

The rest of the paper is organised as follows. Section~\ref{sec:related}
reviews related work. Section~\ref{sec:methodology} describes the data, the
features, the decomposition strategies, the models, and the evaluation
protocol. Section~\ref{sec:results} presents the results.
Section~\ref{sec:discussion} interprets them and states their limitations, and
Section~\ref{sec:conclusion} concludes.

% ---------------------------------------------------------------------------
\section{Related Work}\label{sec:related}
% ---------------------------------------------------------------------------

Three lines of research meet in this study: decomposition-based forecasting,
global forecasting models, and feature-based meta-learning.

Decomposition has a long history as a preprocessing step. Zhang and
Qi~\cite{zhang2005neural} showed that neural networks struggle on raw seasonal
series and improve substantially after deseasonalisation and detrending.
Theodosiou~\cite{theodosiou2011stl} found decomposition-based forecasting
competitive on monthly and quarterly M-competition data, the same setting we
study. Bergmeir et al.~\cite{bergmeir2016bagging} used STL to build bootstrap
ensembles for exponential smoothing, exploiting the decomposition for variance
reduction rather than direct forecasting. The idea has since moved inside
model architectures. Autoformer~\cite{wu2021autoformer} interleaves
decomposition blocks with attention, and Zeng et
al.~\cite{zeng2023dlinear} pair a simple trend and remainder split with linear
layers and outperform far larger transformers on long-horizon benchmarks. Our
study differs from all of these in treating the decomposition decision itself
as the object of analysis rather than fixing it in advance.

The second line concerns global models, which train a single model across many
series. Smyl~\cite{smyl2020hybrid} won the M4 competition with a hybrid that
learns globally while deseasonalising each series individually. Hewamalage et
al.~\cite{hewamalage2021rnn} review recurrent global models and note that
their behaviour depends strongly on how homogeneous the training pool is.
Bandara et al.~\cite{bandara2020clustering} made that observation operational
by training recurrent networks on clusters of similar series and reported
consistent gains. Our bucket-specialised training follows the same intuition,
but the groups are defined by interpretable features rather than learned
clusters, which makes them reproducible and easy to describe.

The third line is meta-learning for forecasting. Wang et
al.~\cite{wang2009rules} induced rules that link series characteristics to the
best forecasting method. Kang et al.~\cite{spectral_entropy_ref} placed the M3
series in a feature-based instance space to expose where methods succeed and
fail. Talagala et al.~\cite{talagala2023metalearning} trained a classifier
that selects a method from series features, and
FFORMA~\cite{montero2020fforma} extended the idea to feature-based weighting
of a model pool, finishing second in M4. All of this work selects or combines
models. The preprocessing decision, decompose or not, is made silently before
model selection even starts, and to our knowledge it has not been studied as a
feature-conditioned decision in its own right.

We deliberately restrict the study to STL. It is the most widely deployed
decomposition method, its robust variant handles outliers, and its additive
components recompose by simple addition, which keeps the experimental design
clean. The conclusions therefore concern STL specifically.
Section~\ref{sec:discussion} discusses what would be needed to extend them to
other decomposition families.

% ---------------------------------------------------------------------------
\section{Methodology}\label{sec:methodology}
% ---------------------------------------------------------------------------

Our pipeline has five stages: sample series from the benchmark pools, compute
features, decompose, forecast under every combination of decomposition
strategy and model, and evaluate everything under a common protocol. This
section describes each stage in turn.

\subsection{Problem Formulation}\label{sec:problem}

Let $\mathbf{y} = (y_1, y_2, \ldots, y_T)$ be a univariate time series with
seasonal period $m$ and forecast horizon $h$. Given a set of forecasting
methods $\mathcal{M}$ and a set of decomposition strategies $\mathcal{D}$, we
ask whether observable properties of $\mathbf{y}$ can predict which pair
$(d, M) \in \mathcal{D} \times \mathcal{M}$ minimises forecast error for that
series.

To study this, we describe each series by a six-dimensional feature vector
$\phi(\mathbf{y})$, partition the series into tercile groups along each
feature dimension, train every combination of strategy and model within each
group, and compare performance across groups. If decomposition helps only
under certain conditions, the signal should appear as a systematic performance
difference between feature terciles.

\subsection{Datasets}\label{sec:data}

We use two established benchmarks from the M-Competition series, the
M3~\cite{makridakis2000m3} and M4~\cite{makridakis2020m4} datasets, restricted
to monthly and quarterly frequencies. Table~\ref{tab:datasets} summarises the
collection.

\begin{table}[ht]
\caption{Dataset summary. Horizon and season length follow the competition
definitions. Series counts reflect the representative sample described in
Section~\ref{sec:sampling}.}\label{tab:datasets}
\begin{tabular}{llccc}
\toprule
\textbf{Dataset} & \textbf{Frequency} & \textbf{Season $m$} &
\textbf{Horizon $h$} & \textbf{Series} \\
\midrule
M3 Monthly    & Monthly   & 12 & 18 & 750 \\
M4 Monthly    & Monthly   & 12 & 18 & 750 \\
M3 Quarterly  & Quarterly & 4  & 8  & 750 \\
M4 Quarterly  & Quarterly & 4  & 8  & 750 \\
\midrule
\multicolumn{4}{l}{Total} & 3{,}000 \\
\botrule
\end{tabular}
\end{table}

\subsubsection{Representative Sampling.}\label{sec:sampling}

The M3 and M4 pools contain several thousand series per frequency, so we draw
750 series per source dataset with stratified quantile sampling. Series are
ranked by each feature $\phi_k$ and divided into $S = 30$ strata of equal
size. The quota of 750 series is distributed across strata proportionally and
filled by draws without replacement, so the sample reproduces the marginal
feature distribution of the full pool. A fixed random seed makes the sampling
deterministic.

\subsection{Time Series Feature Extraction}\label{sec:features}

Six features characterise each series. Where noted, they operate on the STL
residuals described in Section~\ref{sec:stl}, so that they capture intrinsic
behaviour rather than removable seasonal structure. All features are computed
from training observations only, with no access to the test period. Four of
the six features are computed from STL components, which ties the features to
the decomposition under study. Section~\ref{sec:discussion} returns to this
point.

\paragraph{Non-normality, $\phi_1$.}
Residuals $r_t$ from the STL fit are standardised as
$z_t = (r_t - \tilde{r}) / (1.4826 \cdot \mathrm{MAD}(r) + \varepsilon)$,
where $\tilde{r}$ is the median. We apply the D'Agostino omnibus normality
test~\cite{dagostino1971} to $\{z_t\}$ and convert its $p$ value to a
log-scale score:
\begin{equation}
  \phi_1 = \min\!\bigl(-\log_{10}(p_1 + 10^{-12}),\; 12\bigr).
  \label{eq:phi1}
\end{equation}
High values indicate strong departure from Gaussian residuals.

\paragraph{Nonlinearity, $\phi_2$.}
We test whether the lagged history contains nonlinear predictive structure by
comparing a linear and a nonlinear autoregressive model. Let
$\ell = \min(m, 12)$ be the embedding dimension. We build lag-feature matrices
from $\{(y_{t-\ell},\ldots,y_{t-1}, y_t)\}$ and split the data chronologically
at 55, 70, and 85 percent of the series length to obtain three train and test
folds. On each fold we fit a ridge regression~\cite{hoerl1970ridge} with
$\alpha = 1$ and a random forest~\cite{breiman2001random} with 40 trees of
maximum depth four. The feature is the mean relative improvement of the
nonlinear model:
\begin{equation}
  \phi_2 = \min\!\left(
    \max\!\left(0,\;
    \frac{\overline{\mathrm{MAE}}_{\mathrm{linear}} -
           \overline{\mathrm{MAE}}_{\mathrm{RF}}}
         {\overline{\mathrm{MAE}}_{\mathrm{linear}} + 10^{-8}}
    \right),\; 2
  \right),
  \label{eq:phi2}
\end{equation}
where $\overline{\mathrm{MAE}}$ is the mean across folds. In plain terms,
$\phi_2$ is the relative error reduction of the random forest over ridge
regression, floored at zero and capped at two. A value of zero means the
random forest offers no improvement, which is structurally common in linear
datasets, as Section~\ref{sec:nonlin_edge} shows.

\paragraph{Spectral entropy, $\phi_3$.}
We compute the power spectral density of the demeaned series via a periodogram
and derive its normalised Shannon entropy:
\begin{equation}
  \phi_3 = -\frac{\displaystyle\sum_{k} p_k \ln p_k}{\ln K},
  \label{eq:phi3}
\end{equation}
where $p_k = \hat{S}(f_k) / \sum_j \hat{S}(f_j)$ are the normalised spectral
weights and $K$ is the number of non-zero frequencies. The result lies in
$[0, 1]$. Low values indicate a spectrum concentrated at a few dominant
frequencies, which means strong and regular periodicity. High values indicate
a diffuse, unpredictable spectrum~\cite{spectral_entropy_ref}.

\paragraph{Evolving seasonality, $\phi_4$.}
To detect seasonal behaviour that changes over time, we compute the seasonal
strength of Equation~\ref{eq:seasstrength} in rolling windows of length
$\max(36, 3m)$ for monthly and $\max(16, 4m)$ for quarterly series, advancing
by $\lfloor m/2 \rfloor$ observations per step. The seasonal strength in each
window is
\begin{equation}
  F_s^{(w)} = \max\!\left(0,\; 1 - \frac{\mathrm{Var}(R^{(w)})}
                                         {\mathrm{Var}(S^{(w)} + R^{(w)})}\right),
  \label{eq:seasstrength}
\end{equation}
where $S^{(w)}$ and $R^{(w)}$ are the STL seasonal and residual components in
window $w$. The feature is the standard deviation of the sequence
$\{F_s^{(w)}\}$. A high value indicates that seasonal strength changes
substantially over time.

\paragraph{Structural break strength, $\phi_5$.}
We detect level shifts in the nonseasonal component $T_t + R_t$ with the PELT
algorithm~\cite{killick2012}, using an RBF cost function and a penalty
proportional to $\log(n)$. Let $\mathcal{B}$ be the set of breakpoints
detected in the central 60 percent of the series, which avoids boundary
effects. For each candidate break $b \in \mathcal{B}$, the mean-shift contrast
is
\begin{equation}
  \Delta_b = \frac{|\bar{x}_{<b} - \bar{x}_{\geq b}|}
                  {1.4826 \cdot \mathrm{MAD}(T + R) + \varepsilon},
  \label{eq:breakcontrast}
\end{equation}
and $\phi_5 = \min(\max_b \Delta_b,\; 20)$. The feature measures the size of
the largest structural shift relative to the scale of the nonseasonal
component.

\paragraph{Conditional heteroscedasticity, $\phi_6$.}
We test for ARCH effects~\cite{engle1982arch} in the normalised STL residuals
$\{z_t\}$ using Engle's LM test with lag order $\min(12, \lfloor n/5 \rfloor)$
and return
\begin{equation}
  \phi_6 = \min\!\bigl(-\log_{10}(p_6 + 10^{-12}),\; 12\bigr),
  \label{eq:phi6}
\end{equation}
where $p_6$ is the test $p$ value. High values indicate serially correlated
residual variance, a pattern known as volatility clustering.

\subsubsection{Feature Bucket Assignment.}\label{sec:buckets}

For each feature $\phi_k$ and frequency, all selected series are ranked by
feature value and assigned to one of three buckets of equal size. Low contains
the bottom third, Medium the middle third, and High the top third. The bucket
boundaries are the $33.3^{\text{rd}}$ and $66.7^{\text{th}}$ percentiles
computed across the combined M3 and M4 series. This partitioning drives the
bucket-specialised training described in Section~\ref{sec:finetuning}.

\subsection{STL Decomposition Framework}\label{sec:stl}

Seasonal and trend decomposition using Loess, hereafter
STL~\cite{cleveland1990stl}, factorises a time series as
\begin{equation}
  y_t = T_t + S_t + R_t,
  \label{eq:stl}
\end{equation}
where $T_t$ is a locally smooth trend, $S_t$ is a periodic seasonal component
that is approximately zero-mean within each period, and $R_t$ is the
remainder. We use the robust variant of STL, which down-weights outlier
observations during the Loess smoothing passes, as implemented in
\texttt{statsmodels}~\cite{seabold2010}.

We evaluate three strategies for using this decomposition.

\begin{description}
  \item[\textbf{Direct}.] The forecasting model operates on the raw series
    $\mathbf{y}$ and no decomposition is applied.

  \item[\textbf{STL Seasonal Naive}.] The model is trained on the nonseasonal
    component $N_t = T_t + R_t$ and produces a forecast $\hat{N}_{t+1:t+h}$.
    The seasonal component is extrapolated by repeating the last observed full
    seasonal cycle, $\hat{S}_{t+j} = S_{t+j-m}$ for $j = 1, \ldots, h$. The
    final forecast is $\hat{y}_{t+j} = \hat{N}_{t+j} + \hat{S}_{t+j}$.

  \item[\textbf{STL All Components}.] Three separate model instances are
    trained on $T_t$, $S_t$, and $R_t$, yielding component forecasts
    $\hat{T}$, $\hat{S}$, and $\hat{R}$. The final forecast is the sum:
    \begin{equation}
      \hat{y}_{t+j} = \hat{T}_{t+j} + \hat{S}_{t+j} + \hat{R}_{t+j},
      \quad j = 1, \ldots, h.
      \label{eq:recompose}
    \end{equation}
\end{description}

Tables abbreviate the three strategies as Direct, STL-SN, and STL-AC.

\subsection{Forecasting Model Families}\label{sec:models}

We benchmark ten models from four families across two training paradigms.
Statistical models fit one series at a time. All other families train global
models on the full cohort of series.

\paragraph{Statistical.}
AutoARIMA~\cite{hyndman2008automatic}, AutoSARIMA with seasonal differencing
forced to $D = 1$, and AutoETS~\cite{hyndman2008exponential}, all implemented
via \texttt{statsforecast}~\cite{olivares2022statsforecast}.

\paragraph{Machine learning.}
AutoRidge, AutoLinearRegression, AutoXGBoost~\cite{chen2016xgboost}, and
AutoRandomForest, trained jointly on all cohort series with sliding-window lag
features via \texttt{MLForecast}~\cite{garza2022mlforecast}. Hyperparameters
are tuned with Optuna~\cite{akiba2019optuna} over 100 trials.

\paragraph{Neural.}
AutoNLinear, AutoNHITS~\cite{challu2023nhits}, and AutoLSTM, implemented via
\texttt{NeuralForecast}~\cite{olivares2023neuralforecast}.

\paragraph{Transformers.}
AutoTFT~\cite{lim2021tft} and AutoPatchTST~\cite{nie2022patchtst}, also
implemented via \texttt{NeuralForecast}.

\subsection{Bucket-Specialised Training}\label{sec:finetuning}

For the global model families we investigate whether regime-specialised models
outperform a single global model trained on all 750 series of a cohort. The
hypothesis is that series sharing the same feature tercile have sufficiently
similar dynamics to reward a dedicated model.

Each task configuration combines one feature, one frequency, one decomposition
strategy, and one model. For each configuration we train three independent
global models, one per tercile bucket. Each model trains exclusively on the
roughly 250 series in its bucket and is evaluated exclusively on test series
from the same bucket. The single global model, trained on all 750 series with
no bucket restriction, serves as the comparison baseline throughout.
Performance is measured with the specialisation delta $\Delta_{\mathrm{FT}}$
defined in Section~\ref{sec:evaluation}, where a positive value means the
specialised model wins.

Statistical models are excluded because they fit each series independently, so
the composition of the training pool cannot affect them. Every global
configuration is therefore evaluated twice, once with all 750 series as the
global baseline and once per tercile bucket under the specialised condition.

\subsection{Evaluation Protocol}\label{sec:evaluation}

All experiments share the evaluation protocol described in this section.

\subsubsection{Cross-Validation.}
We use rolling-origin cross-validation~\cite{tashman2000} with $W = 3$ windows
per series, advancing by $h$ observations per window. For each series $i$ and
window $w$, the training set ends at a cutoff $c_i^{(w)}$ and the test set
contains exactly the next $h$ observations.

\subsubsection{Metric.}
The primary metric is the relative naive error, RelNaive for short:
\begin{equation}
  \mathrm{RelNaive} =
    \frac{\mathrm{MAE}(\mathbf{y},\,\hat{\mathbf{y}})}
         {\mathrm{MAE}(\mathbf{y},\,\hat{\mathbf{y}}^{\mathrm{naive}})},
  \label{eq:relnai}
\end{equation}
where $\hat{\mathbf{y}}^{\mathrm{naive}}$ is the seasonal naive forecast that
repeats the last observed seasonal cycle,
$\hat{y}_{t+j}^{\mathrm{naive}} = y_{t+j-m}$ for $j = 1, \ldots, h$. A value
below 1 means the model beats the seasonal naive baseline. We also report a
clipped variant $\mathrm{RelNaive}^{(10)}$ that caps values at 10, which
limits the influence of catastrophic failures on aggregate rankings. Aggregate
statistics are means over all series and windows within each reporting group,
where a group fixes the feature, the decomposition method, the model family,
and the training scope.

The STL delta measures the improvement from applying STL decomposition
relative to the Direct strategy:
\begin{equation}
  \Delta_{\mathrm{STL}} =
    \mathrm{RelNaive}^{(10)}_{\mathrm{Direct}} -
    \mathrm{RelNaive}^{(10)}_{\mathrm{STL}},
  \label{eq:stl_delta}
\end{equation}
so a positive value means STL reduces error. The specialisation delta measures
the improvement of a bucket-specialised model over the global baseline:
\begin{equation}
  \Delta_{\mathrm{FT}} =
    \mathrm{RelNaive}^{(10)}_{\mathrm{global}} -
    \mathrm{RelNaive}^{(10)}_{\mathrm{bucket}},
  \label{eq:ft_delta}
\end{equation}
where again a positive value favours the specialised model.

We assess the statistical significance of aggregate deltas with the Wilcoxon
signed-rank test~\cite{wilcoxon1945} applied at the series level, correcting
$p$ values for multiple comparisons with Holm's method~\cite{holm1979}.

% ---------------------------------------------------------------------------
\section{Results}\label{sec:results}
% ---------------------------------------------------------------------------

We report results in five parts: the global baselines, the effect of STL
decomposition, the effect of bucket-specialised training, an edge case in the
feature bucketing, and the statistical significance of all comparisons.
Section~\ref{sec:summary} collects the main findings.

% ---------------------------------------------------------------------------
\subsection{Global Baseline Performance}\label{sec:baseline}
% ---------------------------------------------------------------------------

All 3{,}000 series, 1{,}500 monthly and 1{,}500 quarterly, are evaluated over
three rolling cross-validation windows. Table~\ref{tab:top10} ranks the global
model configurations by mean $\mathrm{RelNaive}^{(10)}$ over all series with
no bucket restriction. AutoETS without decomposition is the best single
configuration, with a mean of 0.888 and a median of 0.791, and the only
configuration without preprocessing whose mean beats the naive forecast.
STL-SN paired with statistical or transformer models dominates the next nine
places.

\begin{table}[ht]
\caption{Top 10 global model configurations by mean $\mathrm{RelNaive}^{(10)}$,
where lower is better. Direct means no decomposition, STL-SN is STL Seasonal
Naive, and STL-AC is STL All Components.}\label{tab:top10}
\begin{tabular*}{\textwidth}{@{\extracolsep\fill}llllcc}
\toprule
Rank & Family & Model & Decomp.\ & Mean & Median \\
\midrule
1  & Statistical  & AutoETS               & Direct  & 0.888 & 0.791 \\
2  & Statistical  & AutoARIMA             & STL-SN  & 0.930 & 0.887 \\
3  & Neural       & AutoNLinear           & STL-SN  & 0.931 & 0.889 \\
4  & Statistical  & AutoSARIMA            & STL-SN  & 0.931 & 0.889 \\
5  & ML           & AutoLinearRegression  & STL-SN  & 0.933 & 0.893 \\
6  & ML           & AutoRidge             & STL-SN  & 0.933 & 0.889 \\
7  & Transformer  & AutoPatchTST          & STL-SN  & 0.933 & 0.891 \\
8  & Statistical  & AutoETS               & STL-SN  & 0.935 & 0.918 \\
9  & Transformer  & AutoPatchTST          & STL-AC  & 0.936 & 0.826 \\
10 & Transformer  & AutoTFT               & STL-SN  & 0.936 & 0.919 \\
\botrule
\end{tabular*}
\end{table}

A persistent divergence between mean and median runs through the whole study.
Under STL decomposition the ML and Neural families have mean
$\mathrm{RelNaive}^{(10)}$ between 1.12 and 1.19 while their medians stay near
or below 1. A heavy right tail of catastrophic failures on a minority of
series pulls the mean above the naive baseline even when most series benefit.
We therefore report both statistics throughout.

% ---------------------------------------------------------------------------
\subsection{Effect of STL Decomposition}\label{sec:stl_results}
% ---------------------------------------------------------------------------

Table~\ref{tab:stl_delta} reports the STL delta for the global model. Positive
values mean STL reduces error relative to Direct.

\begin{table}[ht]
\caption{STL delta $\Delta_{\mathrm{STL}}$ by decomposition and model family
for the global model. Mean and median are taken over all series and windows.
The win rate is the fraction of series where STL helps.}\label{tab:stl_delta}
\begin{tabular*}{\textwidth}{@{\extracolsep\fill}llccc}
\toprule
Decomposition & Family & Mean $\Delta$ & Median $\Delta$ & Win Rate \\
\midrule
STL-AC & Transformer  & $+0.039$ & $+0.141$ & $68.5\%$ \\
STL-SN & Transformer  & $+0.063$ & $+0.091$ & $67.9\%$ \\
STL-AC & Neural       & $-0.119$ & $+0.063$ & $57.6\%$ \\
STL-SN & Neural       & $-0.109$ & $+0.027$ & $55.4\%$ \\
STL-AC & ML           & $-0.191$ & $+0.023$ & $52.5\%$ \\
STL-SN & ML           & $-0.194$ & $-0.002$ & $49.5\%$ \\
STL-AC & Statistical  & $-0.028$ & $-0.085$ & $39.8\%$ \\
STL-SN & Statistical  & $-0.009$ & $-0.109$ & $37.1\%$ \\
\botrule
\end{tabular*}
\end{table}

\paragraph{Transformers.}
Transformers benefit robustly from STL. Both STL-SN and STL-AC give positive
mean and median deltas, with win rates near 68\%. The improvement is strongest
under STL-AC, with a median of $+0.141$, which suggests that training separate
global submodels for trend, seasonal, and residual components fits the
inductive biases of PatchTST and TFT well.

\paragraph{Neural and ML models.}
These families show the divergence between mean and median at its most
extreme. For Neural the median is positive, $+0.063$ under STL-AC and $+0.027$
under STL-SN, but the mean is deeply negative at $-0.119$ and $-0.109$. STL
helps most series but fails badly on series with non-stationary seasonal
patterns, where the decomposition itself is unreliable. ML fares worse, with
an STL-SN mean of $-0.194$ and a median of essentially zero.

\paragraph{Statistical models.}
Both decomposition strategies harm statistical models, with win rates between
37\% and 40\% and negative means and medians. AutoARIMA, AutoSARIMA, and
AutoETS already model trend and seasonality internally. STL preprocessing adds
a redundant transformation and extra parameter uncertainty, so accuracy
degrades rather than improves.

Frequency effects are marginal. The median STL-AC gain is $+0.033$ for monthly
series and $+0.024$ for quarterly series, with no directional difference
between the two frequencies.

% ---------------------------------------------------------------------------
\subsection{Bucket-Specialised Training}\label{sec:ft_results}
% ---------------------------------------------------------------------------

This section measures whether a global model trained only on the series of one
feature bucket beats the single global model trained on all series.

\subsubsection{Overall Results}\label{sec:ft_overall}

Across 1{,}449{,}738 series-window evaluations, the specialised models achieve
a median $\Delta_{\mathrm{FT}}$ of $+0.050$, a mean of $-0.029$, and a win
rate of 56.1\%. The gap between mean and median mirrors the STL results.
Specialisation benefits the majority of series but fails badly on a tail, and
that tail drags the mean negative.

Medium-bucket series improve most consistently, with a median of $+0.059$ and
a win rate of 56.8\%. Low and High are comparable, with medians of $+0.048$
and $+0.046$. The Medium advantage is intuitive. These are the modal series of
each cohort, so the specialised model is not pulled toward the extremes
present in either tail.

\subsubsection{Non-Linear Versus Linear Models}\label{sec:linvnonlin}

The clearest split in the study separates non-linear from linear models.
Table~\ref{tab:lin_nonlin} aggregates over everything else.

\begin{table}[ht]
\caption{Specialisation delta by model type. Results aggregate over all
families, buckets, decompositions, and series.}\label{tab:lin_nonlin}
\begin{tabular}{lccc}
\toprule
Model type & Mean $\Delta_{\mathrm{FT}}$ & Median $\Delta_{\mathrm{FT}}$ & Win rate \\
\midrule
Non-linear & $+0.082$ & $+0.092$ & $60.1\%$ \\
Linear     & $-0.251$ & $-0.014$ & $48.1\%$ \\
\botrule
\end{tabular}
\end{table}

Non-linear models benefit. Focusing training on roughly 250 regime-homogeneous
series improves generalisation within that regime, and both mean and median
are positive, so the gain is robust. Linear models are harmed. Roughly 250
training series are too few for stable coefficient estimation, and the
distribution shift induced by bucketing makes the problem worse. The win rate
of 48.1\% sits below chance, which signals a directional harm rather than
noise. Section~\ref{sec:significance} confirms the significance.

Table~\ref{tab:per_model} gives the per-model breakdown.

\begin{table}[ht]
\caption{Specialisation delta per model, aggregated over all buckets,
decompositions, features, and frequencies. Linear models are marked L and
non-linear models NL.}\label{tab:per_model}
\begin{tabular*}{\textwidth}{@{\extracolsep\fill}lllccc}
\toprule
Family & Model & Type & Mean $\Delta$ & Median $\Delta$ & Win rate \\
\midrule
ML          & AutoXGBoost           & NL & $+0.236$ & $+0.187$ & $65.3\%$ \\
Neural      & AutoLSTM              & NL & $+0.231$ & $+0.162$ & $64.2\%$ \\
ML          & AutoRandomForest      & NL & $+0.146$ & $+0.155$ & $61.6\%$ \\
Neural      & AutoNHITS             & NL & $-0.031$ & $+0.076$ & $59.2\%$ \\
Transformer & AutoTFT               & NL & $-0.076$ & $+0.059$ & $57.5\%$ \\
Transformer & AutoPatchTST          & NL & $-0.013$ & $+0.014$ & $52.6\%$ \\
\midrule
ML          & AutoRidge             & L  & $-0.327$ & $-0.004$ & $49.5\%$ \\
ML          & AutoLinearRegression  & L  & $-0.327$ & $-0.002$ & $49.7\%$ \\
Neural      & AutoNLinear           & L  & $-0.099$ & $-0.029$ & $45.1\%$ \\
\botrule
\end{tabular*}
\end{table}

AutoXGBoost and AutoLSTM gain the most. AutoXGBoost reaches a median of
$+0.187$ with a win rate of 65.3\%, and AutoLSTM a median of $+0.162$ with a
win rate of 64.2\%. Both benefit from the focused within-bucket training
distribution. AutoNHITS and the transformers have positive medians but
negative means, the signature of fat-tailed error distributions. AutoRidge and
AutoLinearRegression suffer the largest harm, with a mean delta of $-0.327$, a
direct consequence of coefficient instability on small training pools. The
practical rule is simple: restrict bucket specialisation to non-linear models.

\subsubsection{Interaction with STL Decomposition}\label{sec:ft_stl_interact}

Specialisation also interacts with the decomposition strategy, as
Table~\ref{tab:ft_decomp} shows.

\begin{table}[ht]
\caption{Specialisation delta by decomposition strategy. All non-statistical
model families and buckets are combined.}\label{tab:ft_decomp}
\begin{tabular}{lccc}
\toprule
Decomposition & Mean $\Delta_{\mathrm{FT}}$ & Median $\Delta_{\mathrm{FT}}$ & Win rate \\
\midrule
Direct   & $-0.156$ & $+0.108$ & $60.4\%$ \\
STL-SN   & $+0.044$ & $+0.032$ & $55.6\%$ \\
STL-AC   & $+0.026$ & $+0.021$ & $52.2\%$ \\
\botrule
\end{tabular}
\end{table}

Direct gives the highest median, $+0.108$, but the worst mean, $-0.156$.
Without seasonal preprocessing, a model trained on roughly 250 series overfits
the idiosyncratic structure of its bucket. STL-SN is the most robust pairing,
with a mean of $+0.044$ and a median of $+0.032$. Removing seasonality first
gives the model a more stationary target and prevents most catastrophic
failures.

% ---------------------------------------------------------------------------
\subsection{The Nonlinearity Edge Case}\label{sec:nonlin_edge}
% ---------------------------------------------------------------------------

The nonlinearity feature $\phi_2$ exposes a structural limitation of tercile
bucketing. Because $\phi_2$ is floored at zero in Equation~\ref{eq:phi2}, the
M3 and M4 series accumulate at exactly zero: 81.8\% of monthly and 92.8\% of
quarterly series. Both tercile boundaries therefore coincide at
$q_{1/3} = q_{2/3} = 0$. Every zero-valued series falls into Low, every
non-zero series falls into High, and the Medium bucket is empty.
Table~\ref{tab:nonlin_buckets} shows the resulting bucket sizes.

\begin{table}[ht]
\caption{Bucket sizes for the nonlinearity feature $\phi_2$. The coincident
boundaries $q_{1/3} = q_{2/3} = 0$ force a binary split into Low and
High.}\label{tab:nonlin_buckets}
\begin{tabular}{lcccc}
\toprule
Frequency & Low & Medium & High & Total \\
\midrule
Monthly   & 1{,}227 (81.8\%) & \textbf{0} & 273 (18.2\%) & 1{,}500 \\
Quarterly & 1{,}392 (92.8\%) & \textbf{0} & 108 (7.2\%)  & 1{,}500 \\
\botrule
\end{tabular}
\end{table}

The 27 affected task configurations, 1.9\% of all bucket tasks, are excluded
from comparative analyses. The binary split that remains is still meaningful.
Low collects the series where lagged ridge regression is near optimal, and
High collects the minority with genuinely exploitable nonlinear structure.
This is in fact the sharpest regime contrast among all six features.

% ---------------------------------------------------------------------------
\subsection{Statistical Significance}\label{sec:significance}
% ---------------------------------------------------------------------------

We applied Wilcoxon signed-rank tests with Holm correction to every
combination of model, decomposition, and bucket, which gives $n = 81$ tests.
Of these, 55 combinations improve significantly and 23 degrade significantly,
all with $p_{\mathrm{Holm}} < 0.001$, while 3 show no significant change. In
relative terms, specialisation helps significantly in 67.9\% of configurations
and harms significantly in 28.4\%.

The five most improved combinations all pair a non-linear model with the
Medium bucket and STL decomposition, as Table~\ref{tab:top5_sig} shows. The
combination of STL-SN, the Medium bucket, and AutoXGBoost is the single
strongest result in the study, with a median $\Delta_{\mathrm{FT}}$ of
$+0.334$ and a win rate of 72.3\%.

\begin{table}[ht]
\caption{Top 5 significantly improved combinations of model and bucket. All
pass the Wilcoxon test with Holm correction at
$p_{\mathrm{Holm}} < 0.001$.}\label{tab:top5_sig}
\begin{tabular*}{\textwidth}{@{\extracolsep\fill}llllcc}
\toprule
Family & Model & Decomp.\ & Bucket & Med.\ $\Delta$ & Win rate \\
\midrule
ML          & AutoXGBoost      & STL-SN & Medium & $+0.334$ & $72.3\%$ \\
ML          & AutoXGBoost      & STL-AC & Medium & $+0.308$ & $66.9\%$ \\
Neural      & AutoLSTM         & STL-SN & Medium & $+0.299$ & $69.2\%$ \\
ML          & AutoRandomForest & STL-AC & Medium & $+0.291$ & $64.4\%$ \\
ML          & AutoRandomForest & STL-SN & Medium & $+0.290$ & $68.2\%$ \\
\botrule
\end{tabular*}
\end{table}

The significantly harmed combinations are exclusively the linear models and
AutoNLinear. The worst case is AutoNLinear under STL-AC, with a median
$\Delta_{\mathrm{FT}}$ of $-0.088$ and a win rate of 34.5\%. Linear models
under STL-AC suffer twice over. The global model must fit trend, seasonal, and
residual components of roughly 250 series at once, which leaves each linear
submodel too few degrees of freedom while recomposition compounds the errors.

% ---------------------------------------------------------------------------
\subsection{Summary of Findings}\label{sec:summary}
% ---------------------------------------------------------------------------

Four findings summarise the study.

\begin{enumerate}
  \item STL decomposition effects are family specific, as
    Section~\ref{sec:stl_results} shows. Transformers benefit robustly, with
    positive mean and median deltas and win rates near 68\%. Neural and ML
    models show positive medians but severely negative means, because STL
    fails on non-stationary series. Statistical models are uniformly harmed,
    since their internal trend and seasonal modelling makes STL preprocessing
    redundant.

  \item Bucket specialisation improves non-linear models and harms linear
    ones, as Section~\ref{sec:linvnonlin} shows. Non-linear architectures gain
    a median of $+0.092$ with a win rate of 60.1\%, while linear models lose a
    median of $-0.014$ with a win rate of 48.1\%. The split is statistically
    significant in 67.9\% of tested configurations.

  \item The Medium bucket is the most reliable and STL-SN is the safest
    pairing, as Sections~\ref{sec:ft_overall} and \ref{sec:ft_stl_interact}
    show. Medium-bucket series yield the most consistent specialisation gains,
    and applying STL-SN before bucket training gives positive mean and median
    deltas, eliminating the catastrophic failures seen under Direct.

  \item Tercile bucketing collapses to a binary split when a feature has a
    degenerate distribution, as Section~\ref{sec:nonlin_edge} shows. The zero
    floor of $\phi_2$ empties the Medium bucket for 1.9\% of tasks. Tercile
    partitioning needs a minimum-support guard.
\end{enumerate}

% ---------------------------------------------------------------------------
\section{Discussion}\label{sec:discussion}
% ---------------------------------------------------------------------------

The central message of the study is that the value of seasonal decomposition
depends far more on the forecasting model than is usually acknowledged.
Statistical methods lose accuracy under STL preprocessing because they already
estimate trend and seasonal structure internally, so the external
decomposition adds a redundant transformation and a second source of
estimation error. Transformers gain because the decomposed components match
their inductive biases, and training one global submodel per component appears
to act as a useful division of labour. Machine learning and neural models sit
in between. They benefit on the majority of series, but they depend on the
decomposition being right, and when seasonality drifts over time the seasonal
naive extrapolation breaks and the recomposed forecast can fail badly.

The same asymmetry explains the persistent gap between means and medians.
Decomposition and specialisation both improve typical performance while
occasionally producing catastrophic failures, and a handful of failures is
enough to reverse a ranking based on means. Any study or production system
that evaluates preprocessing decisions on averages alone risks drawing the
wrong conclusion. We recommend reporting medians and win rates next to means
as a matter of course.

For practitioners the rules that emerge are short. Do not decompose for
statistical models. Decompose for transformers, ideally modelling all
components separately. If you train regime-specialised global models, restrict
them to non-linear architectures, remove seasonality first with the seasonal
naive add-back, and expect the largest gains on series in the middle of the
feature distribution. Treat linear global models trained on small pools with
suspicion.

Several limitations bound these conclusions. The study covers one
decomposition method. STL is the most widely used choice and its additive
components recompose trivially, but other families such as wavelet
decompositions may behave differently and deserve the same treatment. Four of
the six features are computed from STL components, so the features and the
treatment share machinery. A series that STL fits poorly receives both an
unreliable decomposition and an unusual feature value, which makes part of the
observed association self-referential. The study also stops one step short of
full meta-learning, since we never train a selector that maps features to a
recommended strategy. Finally, the evidence comes from monthly and quarterly
M3 and M4 data with pools of roughly 250 series per bucket, and other
frequencies, domains, or pool sizes may shift the boundaries we report.

These limitations point directly at future work. The most natural extensions
are repeating the design with other decomposition methods, training an
explicit meta-learner on the six features, and testing whether pretrained
foundation models for time series respond to deseasonalised inputs the way
transformers trained from scratch do.

% ---------------------------------------------------------------------------
\section{Conclusion}\label{sec:conclusion}
% ---------------------------------------------------------------------------

We asked when seasonal decomposition helps time series forecasting and
answered with a controlled study of 3{,}000 monthly and quarterly M3 and M4
series, six interpretable features, three decomposition strategies, and ten
automatically tuned models. The answer depends mostly on the model family.
Transformers benefit consistently from STL preprocessing, statistical models
are consistently harmed, and machine learning and neural models gain on most
series while failing badly on a minority with unstable seasonality. Training
separate global models per feature bucket helps non-linear architectures and
harms linear ones, and pairing specialisation with seasonal adjustment removes
most of its catastrophic failures.

Beyond the specific numbers, the study makes a broader point. The decision to
decompose is a modelling decision like any other, and it deserves the same
empirical scrutiny that model selection already receives. Series features
computed before any model is fitted offer a cheap and interpretable way to
apply that scrutiny, and our results show they carry real signal about when
preprocessing will pay off.

\backmatter

\bmhead{Acknowledgements}
Not applicable.

\section*{Declarations}

\begin{itemize}
  \item \textbf{Funding:} Not applicable.
  \item \textbf{Conflict of interest:} The author declares no competing interests.
  \item \textbf{Ethics approval:} Not applicable.
  \item \textbf{Data availability:} The M3 and M4 datasets are publicly available
    at \url{https://forecasters.org/resources/time-series-data/}.
  \item \textbf{Code availability:} Source code will be made available upon
    acceptance.
  \item \textbf{Author contribution:} D.C. conceived and designed the study,
    implemented the pipeline, ran all experiments, and wrote the manuscript.
\end{itemize}

\bibliography{references}

\end{document}
